% AY2020 力学 第2問 転記（問題のみ）
\section*{第2問〔II〕}

図2のように質量 $M$，半径 $a$ の円盤の重心 G から距離 $h$ だけ離れた点に軸 S を
取り付け，振り子運動させた。ここで，円盤の振り子運動の角度 $\theta$ は小さく
$\sin\theta\fallingdotseq\theta$ の関係が成り立つ。ここで，重心 G のまわりの円盤の
慣性モーメントを $Ma^2/2$，重力加速度の大きさを $g$，軸 S に働く摩擦は無視できる
ものとする。このとき，次の問い (1) および (2) に答えよ。

\begin{enumerate}[label=(\arabic*)]
  \item 軸 S のまわりの円盤の慣性モーメントを求めよ。
  \item 円盤の振り子運動の周期を求めよ。
\end{enumerate}

\begin{center}
% 図2：重心Gからh離れた点Sを軸とする円盤の物理振り子。鉛直（破線）からθ傾く
\begin{tikzpicture}[>=stealth, scale=1.0]
  \coordinate (S) at (0,1.3);
  % 鉛直破線
  \draw[dashed] (S) -- (0,-1.9);
  % 釣合い位置の円盤（破線、中心は鉛直線上）
  \draw[dashed] (0,0) circle (1.3);
  % 傾いた円盤（実線）。θ=18°、G = S + h(sinθ,-cosθ), h=1.3
  \coordinate (G) at (0.402,0.064);
  \draw (G) circle (1.3);
  \draw (S) -- (G);
  \fill (S) circle (1.6pt); \node[above right] at (S) {$S$};
  \fill (G) circle (1.3pt); \node[right] at (G) {$G$};
  \node at (0.42,0.66) {$h$};
  % θ（鉛直線と SG の間）
  \draw (0,0.75) arc (-90:-72:0.55); \node at (0.18,0.7) {$\theta$};
  \node at (1.5,-1.6) {図 2};
\end{tikzpicture}
\end{center}
